\subsection{Measuring simultaneous coverage}

We measure the spatial extent of humid heat on the 465-site lattice at the
original 121-site peak time. Each retained site represents a
$0.8^\circ\times0.9^\circ$ cell clipped to the study rectangle and the
Natural Earth land polygon. WGS84 cell areas sum to 3.221 million km$^2$,
covering 95.7\% of the mapped land in that rectangle. Let $a_i$ be the
represented land area and $a_+=\sum_i a_i$. At threshold $u$, define
\begin{equation}
 A_t(u)=\frac{\sum_i a_i I(y_{it}>u)}{a_+},\qquad
 C_t(u)=\frac{\max_{B\in\mathcal B_t(u)}\sum_{i\in B}a_i}{a_+},
 \label{eq:coverage}
\end{equation}
where $\mathcal B_t(u)$ contains connected components of exceedance cells
under four-neighbour grid adjacency, with $C_t(u)=0$ when the exceedance set
is empty. Multiplication by $a_+$ gives the corresponding areas in km$^2$.
These quantities describe simultaneous coverage at the sampling resolution.
Spatial extent and contiguity are also central to temperature-event analyses
by \citet{healy2025} and \citet{skinner2025}.

The temperature range uses the 1st and 99.5th percentiles of the two
development summers, rounded outward to 0.25$^\circ$C steps. We retain all
77 thresholds from 9.5 to 28.5$^\circ$C. For each threshold we compare high
and middle days within month-year, then average months and evaluation summers
equally. The differences are expressed in area percentage points and km$^2$.
Pointwise descriptive intervals resample five-year calendar blocks of the
annual contrasts, retaining the positions of the two excluded development
years. Eight-neighbour connectivity and cosine-latitude weights provide
separate sensitivity comparisons.

We then compare how summaries of the 121-site field describe coverage on the
465-site grid. A mean-and-season regression is followed by a strong simple
baseline that also includes spatial variance, the signed latitude slope and
the area-weighted 121-site exceedance proportion. Two further models add
five binned semivariances or the five graph values, each divided by spatial
variance and log-transformed to describe relative scale structure. All four
use the same spline regression with ridge regularisation and fitted
proportions clipped to $[0,1]$. Linear spatial interpolation, followed by
coverage calculation, provides a full-field comparison.

Seven held-out five-year calendar blocks and one-year training buffers
separate evaluation from fitting. Nested blocked validation selects the
ridge penalty using training years, following the dependence-aware approach
to model assessment discussed by \citet{roberts2017}. Errors are averaged
across thresholds, then equally across months and summers, for all days and
for high days separately. An additional exceedance comparison uses only the
344 new sites. The connected-area calculation uses all 465 sites. This
comparison measures spatial information within ERA5-Land across omitted
years. Supplementary Section~S13 gives the complete definitions and fitting
procedure.
